% !TeX root = ../papers/nc-p3.tex
\section{整六次点曲線の全排除と整曲線層の分類}
\label{sec:integral-sextic-exclusion}
\begin{theorem}\label{thm:integral-sextic-exclusion}
標準Koszul左右AS正則窓で四辺の左右が点ファイバーなら、整な被約六次曲線と
これと交わらない零次元スキームが点グラフになることはない。
\end{theorem}
\begin{proof}[位数4と巡回被覆の固有束]
前節により残るのは正規化 $N$ の種数が0または1の整な特異平面四次 $C$ である。
位数4のときは $C\to\mathbb P^1$、$u^4=F_4(v,w)$ の有限平坦巡回被覆を $f$ とする。
$d_i=L_{i+1}/L_i$ は $\sigma^*d_i=d_iK^4=d_i$ を満たし、$d_i^4=1$ である。
実際四共役の積はノルムの引戻しで、次数零なので直線上で自明である。
従って $\sigma^*L_i=L_i d_i^4K^6=L_i$、各非有効次数2束 $T=L_i/\omega_C$ は不変である。
生成元の同型をスカラーで調整して四乗を1にし、$T$ を線型化する。
$f_*T$ の四つの固有層は各々直線束である（一般点で各指標が一回ずつ現れ、直線上で無捩れ）。
$H^0(T)=H^1(T)=0$ なのでこれらは全て $\mathcal O(-1)$ である。
$u$ の乗法は隣接固有束を基底の $\mathcal O(1)$ でひねる四つの一次式となり、積は $F_4$ である。
根の上のファイバーでは $T$ が可逆ゆえ $k[u]/u^4$ の自由階数1加群で、$u$ の作用の階数は3である。
従って同じ根で二つの一次式は消えない。$F_4$ は四単純根を持ち $C$ が滑らかになってしまう。

\smallskip
\noindent\textit{三割線からの二つの自己写像の次数。}
$\tau=\sigma^{-1}$ は恒等または対合とする。前節の三割線の構成を使う。
正規化の種数2を用いたのは、二つの非定数自己写像が得られた\emph{後}のHurwitzの箇所だけである。
従って今も $r,s:N\to N$ は非定数で、一般の滑らかな $p$ に対し
\[
 D_p+r(p)\in |T+s(p)|,\qquad p+\tau p+r(p)\in |U+s(p)|,\qquad
 U=T K .
\]
ここで $D_p$ は長さ2の可動因子である。$U$ は非有効である。
有効なら $U$ と $\omega_C/U$ はそれぞれ高々一つの切断を持ち、一般の $s(p)$ で
$|U+s(p)|$ は固定因子と一点 $s(p)$ の和だけとなる。これは可動な二点 $p,\tau p$ を含めず、
$\tau=1$ のときも一般の $p$ を二回含めない。

非有効な次数2束 $U$ について、一般の滑らかな $x$ に対し $|U+x|$ は一点で、次数3の因子を定める。
これを $N\times N$ に閉じた対応を重複込みで $B_U$ とする。両射影次数は3であり、
転置は $U'=\omega_C/U$ から得る対応である。実際、滑らかな点対では
$H^0(U(x-y))\ne0$ と $H^0(U'(y-x))\ne0$ がSerre双対で同値である。
因子の重複も一致する。一般の点対を動かす局所環上で $R\Gamma(U(x-y))$ は
$H^0(U(x))\to U(x)|_y$ の二直線束の評価で表せ、その行列式が因子を切る。
相対Serre双対は転置側をその双対二項複体と同定するので消滅位数も同じである。
どの特異点 $z$ も一般の $|U+x|$ に入らない。
$\mathfrak m_z U$ は $h^0=0,\chi=-1$ で、そのSerre双対にはただ一つの独立切断があり、
一般点 $x$ で消えない。従って $h^0(\mathfrak m_z U(x))=0$ である。$U'$ 側でも同じで、
両方向で三つの点が一般の滑らかな部分で数えられる。

$m=\deg s$ とすれば、一般の $y$ が $|U+s(p)|$ に入る $p$ を両表示から数えて
$3m=2+\deg r$ を得る。$B_U$ は $(s(p),p)$ と $(s(p),\tau p)$ の像を含む。
異なる像なら $2m\le3$、同じで $\tau=1$ なら重複度が2以上なのでやはり $2m\le3$ である。
同じで $\tau\ne1$ なら $s\tau=s$ で $m$ は偶数かつ $m\le3$、従って $m=2$ である。
この場合一般の $|U+x|$ は $s^{-1}(x)$ と残る一点 $f(x)$ の和となり、
$r=fs$ と $\deg r=4$ から $f:N\to N$ は次数2である。
任意の正則双対化微分 $\eta\in H^0(C,\omega_C)$ を $N$ 上の有理微分と見ると
\[
                            f^*\eta+s_*\eta=\eta .
\]
これは動くCartier因子のAbel写像の微分である。滑らかな点で点が動くときの微分と
$\eta$ の対はその点での微分評価であり（局所パラメータの単純極の留数）、
複数の動点なら和、有限被覆の一ファイバーなら各逆枝の和すなわちtraceを取る。
標準平面の対合は非自明なため反不変な非零 $\eta$ があり、$s\tau=s$ より $s_*\eta=0$ である。

ところが種数0,1の滑らかな曲線の次数2の自己写像は非零の有理微分をそのまま保てない。
実際、各極の位数から1を引いた非負部分の合計は引戻しで2倍になる。
従って極があれば単純極だけで、その台は逆像で不変である。種数1の自己写像はエタールだからこれは不可能で、
正則微分に対する倍率も絶対値 $\sqrt2$ である（複素一様化、有限生成体を複素数に埋めればよい）。
射影直線では留数の和から極が少なくとも二つあり、Hurwitzの全分岐次数2から逆像不変な台は高々二点である。
二点を $0,\infty$ とすれば写像は $cz^{\pm2}$、微分は $dz/z$ の倍で、倍率は $\pm2$ である。
従って $m=2$ は消え、$r,s$ はともに自己同型である。

\smallskip
\noindent\textit{対合の反不変微分。}
$A=s^{-1},\ g=rs^{-1}$ とすると $B_U$ は $A,\tau A,g$ の三グラフの和である。
$g=1$ なら $D_p\sim T$ に反する。任意の $\eta\in H^0(\omega_C)$ に上と同じ微分を取ると
\[
             A^*(1+\tau^*)\eta+g^*\eta=\eta .
\]
標準平面の対合の座標で $x=v/w$ と置けば、関数体は
$F=u^4+a(x)u^2+b(x)=0$ で表せる。標準随伴公式による微分 $dx/F_u$ と
$x\,dx/F_u$ はともに反不変である。これは特異点でも平面の双対化束の切断である。
従って $g$ は両方を保ち、比 $x$ と $F_u$ を保つ。
しかし $F_u=2u(2u^2+a)$ と $F_u^2=4u^2(a^2-4b)$ より $x,F_u$ は全関数体を生成する。
判別式 $a^2-4b$ は整性から恒等零でない。よって $g=1$ となり矛盾する。

\smallskip
\noindent\textit{恒等四移動と正規化の分岐。}
$\tau=1$ では定理\ref{thm:fixed-differential}で $\deg M+4=0$、すなわち
正規化微分の零因子 $B$ の次数は $2g(N)+2$ である。
一枝の最小正位数を $e$ とすると $B$ の係数は $e-1$ で、全枝での和は
$\delta=3-g(N)$ を超えない。実際、正規化局所環でその枝の一様化元の第1から第 $e-1$ 乗の係数は
元の局所環からは得られず、別々の枝でこれらは独立な余長を与える。
従って $N=\mathbb P^1$、$\deg B=2$ である。
$B_U$ とその転置の三グラフ表示から、全双対化微分について
\[
      (2A^*+g^*)\eta=\eta,\qquad
      (2(A^{-1})^*+(g^{-1})^*)\eta=\eta .
\]
$A\ne1$ である。$s=1$ なら $2p+r(p)\sim U+p$ から $U$ が有効になる。

$B$ の台の点 $a$ に零点を持つ射影直線座標を $t$ とすると、$dt/t^2$ は $C$ 上の正則双対化微分である。
実際この微分は $a$ 以外で正則で、$C$ 上の任意の局所函数を掛けても
その枝の一次係数が零なので留数が零である。ここで有限正規化の留数表示
（各局所函数との積の枝留数の和が零であること）を用いた。
$e\ge3$ なら $dt/t^3$ にも同じことが成り立つ。
第一の等式の右辺に一つしか極がないため、$A,g$ は両方 $a$ を固定する。
異なる二点に極が出れば消し合えず、両方が同じ別の点に出れば $a$ の極が作れない。
$a$ での接倍率の逆数を $c,d$ とし、二つの等式の二重極係数を比較すると
$2c+d=1,\ 2/c+1/d=1$、よって $c=1,d=-1$ である。
$B$ の台が二点なら $A$ は異なる二点を固定しかつ接倍率1なので恒等になってしまう。
一点なら $e=3$ で、三重極を使うと $2c^2+d^2=1$ を要し矛盾する。これで全ての六次を除いた。
\end{proof}

\begin{corollary}\label{cor:integral-point-curves-covered}
点ファイバーで、一次元部分が整な被約曲線かつ有限部と交わらない標準AS窓の全体は、
滑らかな平面三次、滑らかな楕円四次、二類の滑らかな楕円五次、
および特異点を許す八係数の整七次の四行で尽きる。各行の正則性と添字保持同型の判定は
系\ref{cor:smooth-point-curves-covered}と定理\ref{thm:integral-septic-boundary}の通りである。
四行の次数は相異なるので行を越える同型はない。
\end{corollary}
\begin{proof}
次数上限は7であり、次数4以下、5、6、7をそれぞれ
定理\ref{thm:reduced-cubic-point-curves}、定理\ref{thm:reduced-quartic-point-curves}、
定理\ref{thm:singular-integral-quintic-exclusion}と本節および前々節の七次の定理で合わせる。
\end{proof}
